\manualchapter{Patches and troubleshooting}{sec:operation}
\subsection{Pitch-tracking percussion}
Keep Noise Period at 0, mute Tone 3's audio, and patch pitch to Tone 3 V/OCT.
Use a short envelope on Noise Level and compare the two LFSR settings.
This gives pitched, repeating noise or a broader noise texture using the
same clock. For fixed-rate percussion, select Noise Period 1--3 instead.

For a bass plus percussion patch, take Tone 1 and Noise on separate outputs.
Remember that taking Tone 1 out breaks its forward contribution to later
outputs. If the noise changes unexpectedly with a melody, check whether its
period is at 0 and following Tone 3. If a mix sounds flattened, lower the
individual level settings before increasing downstream gain.

\subsection{Cable channels, outputs and saved patches}
The widest input selects 1--16 polyphonic chip instances. Voice columns are
oscillators within each instance, not separate cable channels. Inputs read
the matching channel; supply matching pitch and level channel counts rather
than expecting a mono envelope to repeat across a polyphonic pitch cable.

Each output adds the preceding output when that earlier socket is unpatched.
The accumulated signal clips at approximately $\pm5$\,V. To hear the complete
mix, use the last output and leave earlier outputs empty. Patching an earlier
output removes that signal from the following mix. Lower source levels if
the sum clips.

Rack saves the panel parameters. Reset initializes the chip; oscillator
phase is not preserved by a patch save. Pitch/FM update each sample, while
level and other register controls update every 16 samples.
